\subsubsection*{Misc}

\begin{proposition}
\cref{eq:warp-fa-bvp} is well defined.

The displacement problem of torsion is a pure Neumann problem of the
Laplace operator:
\begin{equation*}
\left\{
\begin{array}{rl}
\Delta \varphi &= f \quad \text { in } \quad \PlaneManifold \\
\nabla_{\bm{\nu}} \varphi &= g \quad \text { on } \quad \PlaneBoundary \\
\int_\Omega \varphi &= 0
\end{array}
\right.
\end{equation*}
where, in order for the problem to be well defined one must have \parencite{shapiro2017neumann}:
\[
\int_\Gamma g = \int_{\Omega} f
\]
\end{proposition}

\begin{proof}
Define $X := \nabla \WarpFA$.
%
The strong form of the governing PDE is:
$$
\mathcal{L}_X \volG = \WarpSV \cdot \volE
\quad \text{in } \PlaneManifold,
\qquad
\nabla_\nu \WarpFA = 0 \quad \text{on } \PlaneBoundary
$$

We now apply Gauss' theorem in the form \citep{abraham1988manifolds}:

$$
\int_{\PlaneManifold} \operatorname{div}_\mu X \, \mu = \int_{\PlaneBoundary} \iota_X \mu
$$

Or in standard notation:

$$
\int_{\PlaneManifold} \nabla \cdot (n_G \nabla \WarpFA) \, \dd A = \int_{\PlaneBoundary} n_G \nabla \WarpFA \cdot \nu \, \dd s
$$

Now use the PDE $\nabla \cdot (n_G \nabla \WarpFA) = n_E \WarpSV$ so the volume integral becomes:

$$
\int_{\PlaneManifold} n_E \WarpSV \, dA = \int_{\PlaneBoundary} n_G \nabla \WarpFA \cdot \nu \, ds
$$

Apply Boundary Condition; Given $\nabla_\nu \WarpFA = 0$ on \(\PlaneBoundary\) 
Then:

$$
\nabla \WarpFA \cdot \nu = 0 \quad \text{on } \PlaneBoundary
\quad \Rightarrow \quad
n_G \nabla \WarpFA \cdot \nu = 0
$$

Therefore:

$$
\int_{\PlaneBoundary} n_G \nabla \WarpFA \cdot \nu \, ds = 0
$$

and so the divergence theorem yields:

$$
\int_{\PlaneManifold} n_E \WarpSV \, dA = 0
$$
\end{proof}


\subsection{Laplacian}
So the Laplacian becomes:

$$
\operatorname{div}_{G}(\nabla \WarpSV 
+\FrameBasis \times \boldsymbol{r})
=
\frac{1}{\WeightG} \operatorname{div}_{A} \left(\WeightG (\nabla \WarpSV +\FrameBasis \times \boldsymbol{r})\right)
=
\frac{1}{\WeightG} \nabla \cdot\left(\WeightG (\nabla \WarpSV +\FrameBasis \times \boldsymbol{r})\right)
=0
$$

\begin{equation*}
\begin{aligned}
\mathcal{L}_{\nabla \WarpFA} \volG &= \WarpSV \volE  
\quad\text{in }\Omega, \\
( \operatorname{div}_{\volG}\nabla\WarpFA )\volG
  &=\WarpSV\,\volE.
\end{aligned}
\qquad
\begin{aligned}
\nabla_\nu \WarpFA = 0 \quad \text{on } \PlaneBoundary \\
\ii_{\nu}\dd\WarpFA = 0
\quad\text{on }\PlaneBoundary;
\end{aligned}
\end{equation*}


\subsection{Equivalence of Warping Boundary Conditions}

Consider two functions \( \WarpSV \colon \Omega \to \mathbb{R} \) governed by the same differential equation:
\[
\operatorname{div}_\mu \left( \nabla \WarpSV + \FrameBasis \times \PlanePosition \right) = 0 \quad \text{in } \Omega,
\]
but with differing boundary conditions.

\paragraph{Formulation 1 (homogeneous Neumann):}
\begin{equation}
\nabla_{\bm{\nu}} \WarpSV = 0 \quad \text{on } \partial \Omega.
\label{eq:hom-neumann}
\end{equation}

\paragraph{Formulation 2 (nonhomogeneous Neumann):}
\begin{equation}
\nabla_{\bm{\nu}} \WarpSV = -\FrameBasis \times (\PlanePosition - \bm{\SectionLocation}) \cdot \bm{\nu}
\quad \text{on } \partial \Omega,
\label{eq:nonhom-neumann}
\end{equation}
where \( \SectionLocation \in \mathbb{R}^2 \) is a fixed reference point in the plane (e.g., the shear center), and \( \bm{\nu} \) is the outward unit normal.

\paragraph{Equivalence.}
The two formulations are equivalent up to an additive affine function. That is, if \( \WarpSV \) satisfies \cref{eq:nonhom-neumann}, then
\[
\WarpSV'(\PlanePosition) := \WarpSV(\PlanePosition) + \FrameBasis \cdot (\PlanePosition - \bm{\SectionLocation})
\]
satisfies \cref{eq:hom-neumann}. Conversely, given a solution \( \WarpSV' \) to \cref{eq:hom-neumann}, the function
\[
\WarpSV(\PlanePosition) := \WarpSV'(\PlanePosition) - \FrameBasis \cdot (\PlanePosition - \bm{\SectionLocation})
\]
satisfies \cref{eq:nonhom-neumann}.

\paragraph{Interpretation.}
\cref{eq:nonhom-neumann} treats \( \WarpSV \) as a \emph{pure warping function}, with torsional twist represented explicitly in the boundary condition. \cref{eq:hom-neumann} instead absorbs the affine twist into the interior solution, simplifying the boundary data. Both formulations yield the same shear strain field:
\[
\nabla \WarpSV + \FrameBasis \times \PlanePosition.
\]
